\documentclass{article}

\begin{document}

\section{Introduction}
This document aims to outline the requirements for prototype 1 and how it will be
tested.

\section{Requirements}
The motor should be able to produce a continuous force of 4 to 8N at approximately 40 to 60c. The 
design should be responsive allowing for fast changes in direction or velocity. The design should
be able to run using step/dir/enable and also direct CAN-BUS communication. The design must be able
to run in a x, y configuration with either 1 motor on each axis or 1 motor on the x-axis and two on the y-axis.

\section{Constraints}
A 3d printers main supply is 24 V limited though can push upto 15-20A depending on supply model. The magnet size
is 10mm (axially) by 10mm (radially) due to pre-purchased before design cycle and to decrease the solution space.
The wire sizes available is similar having been pre-purchased hence only 0.2, 0.3, 0.4 mm wire diameters are available.
The AS5311 linear encoder due to being on hand is the encoder that will be used on this design.

\section{Testing}
To test the motor, it will be configured in a X-Y gantry this allows for collecting thermal, magnetic and electric data
about how the motor operates. Which will inform simulation modelling decisions will also help with optimization of the motor.

\vspace{0.5cm}
Testing will begin with just measuring the BACK-EMF from back driving the motor to see if the phase alignment is correct. If it is
than each phase should be 120 degrees out of phase with one another. Than the encoder will be tested to see if its working correctly
and that the MCU can keep up with high down rates. After that the 3 phase driver will be connected to the motor and the encoder which
will allow for true testing of the setup using forward driving from electrical to mechanical energy. If it boots up and jogs a bit than
the PID loop for each axis can be tuned before running a test planar g-code to gather mix thermal, electrical and magnetic data.

\end{document}